
The approach of this research is heavily informed by previous exploration into hardware-software co-design for implantable medical devices. The fundamental challenge of respecting the strict $1^\circ$C thermal limit has historically forced designers to rely on highly specialized ASICs. Kassiri et al.'s closed-loop neurostimulator \cite{7982670}, Neuralink's N1 ASIC \cite{Musk2019}, and the NeuroPace RNS System \cite{Heck2014} represent the state-of-the-art implementations of this rigid baseline. These systems achieve unparalleled energy efficiency for their exact medical tasks by baking clinical logic into the physical silicon at the cost of algorithmic flexibility. Unlike these ASICs, HALO is built on the RVV open-source instruction set architecture (ISA). This enables its binaries to pick up improvements from every update of the open-source GCC backend without a tape-out. The Results chapter demonstrates the payoff: simply updating the compiler resulted in reduced collective execution cost across the core arithmetic and logical operations (excluding Octave-mode comparison tests), and in three cases (\texttt{and} dbl$\times$dbl, \texttt{and} cpx$\times$dbl, \texttt{or} cpx$\times$dbl) the auto-vectorized scalar code now beats the hand-written intrinsics outright.

\section{Translation for High-Level Languages}

% NOTE TO PROFESSOR: should i expand this section to go into more detail about these earlier tools?

Source-to-source code generation has been a topic of interest since the 1980s \cite{taylor1982upward}. While this coincides with the release of MATLAB in the late 1970s, MATLAB did not have efficient transpilation tools until the 1990s with the introduction of FALCON by the University of Illinois at Urbana Champaign. DeRose et al. created a translator that was capable of compiling MATLAB into Fortran or C while giving developers the ability to "interactively apply optimizations and transformations to the code" \cite{derose1996falcon}. FALCON followed an abstract syntax tree construction approach similar to the approach of this research, but it suffered in the type inference phase. MathWorks have also been developing various MATLAB-to-C tools since the 1990s, such as \texttt{mcc} and \texttt{mex}, which were later replaced by MATLAB Coder \cite{mathworks_matlab_coder}. MATLAB Coder is highly capable at parsing MATLAB scripts and generating optimized C and C++ code targeted primarily at advanced desktop processors and commercial ARM architectures. Similarly, for researchers utilizing Python---another dynamically typed language widely used in BCI data analysis and real-time processing pipelines---tools like Cython \cite{behnel2011cython} offer similar benefits.

Although both MATLAB Coder and Cython are highly effective on standard architectures, they remain limited in use for this research as neither tool inherently supports the automated generation of targeted RISC-V Vector (RVV) intrinsics required to meet the ultra-low-power constraints of the HALO platform. Thus, the ts-traversal and RISCV-Matrix integration developed in this work under Professor Rajit Manohar's lab is necessary to bridge the semantic gap between high-level BCI algorithm prototyping and the specific vector execution demands of the HALO embedded hardware.


